\documentclass[11pt]{article}
\usepackage[a4paper,margin=1in]{geometry}
\usepackage{amsmath,amssymb,mathtools}
\usepackage{bm}

\title{Particle System Model}
\author{Flux}
\date{\today}

\begin{document}
\maketitle

\section{State and Dynamics}
Let $n$ particles evolve in $\mathbb{R}^2$ with positions
\[
X(t) = \big(x_1(t),\dots,x_n(t)\big), \qquad x_i(t)\in\mathbb{R}^2.
\]
The motion with friction is
\[
\ddot{x}_i + \alpha \dot{x}_i = -f_E(x_i) - f_I(X,x_i), \qquad i=1,\dots,n,
\]
where $\alpha\ge 0$ is the friction coefficient.

\section{External Force (Context-Free)}
The external term depends only on the local position:
\[
f_E(x) : \mathbb{R}^2 \to \mathbb{R}^2.
\]
In this application, $f_E$ is built from user-controlled radial centers and does not depend on the current full particle configuration $X$.

\section{Internal Force (In-Context)}
Internal interactions depend on the whole configuration $X$:
\[
f_I(X,x_i)
=
\frac{1}{n}\sum_{k=1}^{K}\sum_{j=1}^{n}
v_k\, g_{s_k}(x_i-x_j)\,(x_i-x_j).
\]
The Gaussian head kernel is
\[
g_s(r)=\exp\!\left(-\frac{\|r\|^2}{s^2}\right).
\]
The self-interaction terms ($j=i$) are set to zero in implementation.

\subsection*{Unity On/Off Vector Term}
Define
\[
w_{ij}=
\begin{cases}
\dfrac{x_i-x_j}{\|x_i-x_j\|}, & \text{Unity On},\\[0.75em]
x_i-x_j, & \text{Unity Off}.
\end{cases}
\]

\section{Two-Head Case ($K=2$)}
With two heads:
\[
v_1>0 \quad\text{(repulsive)}, \qquad v_2<0 \quad\text{(attractive)}.
\]
Using UI amplitudes
\[
a_R = |v_1|,\qquad a_A = |v_2|,
\]
we write
\[
v_1 = +a_R,\qquad v_2 = -a_A.
\]
With radii $s_R$ (repulsion) and $s_A$ (attraction), the scalar radial profile is
\[
\phi(r)=a_R \exp\!\left(-\frac{r^2}{s_R^2}\right)-a_A \exp\!\left(-\frac{r^2}{s_A^2}\right),
\]
and
\[
f_I(X,x_i)
=
\frac{1}{n}\sum_{j=1}^{n}
\phi\!\big(\|x_i-x_j\|\big)\,(x_i-x_j),
\]
with the $j=i$ contribution removed.

\subsection*{Attention On/Off}
With the $w_{ij}$ notation above, two internal-force modes are used.

\paragraph{Attention Off}
\[
f_i
=
\frac{1}{n}\sum_{k=1}^{K}\sum_{j=1}^{n}
v_k\, g_{s_k}(x_i-x_j)\,w_{ij}.
\]

\paragraph{Attention On}
\[
f_i
=
\sum_{k=1}^{K}\sum_{j=1}^{n}
v_k\,
\frac{g_{s_k}(x_i-x_j)}{\sum_{\ell=1}^{n} g_{s_k}(x_i-x_\ell)}
\;w_{ij}.
\]

\end{document}
